\documentclass[UTF8,11pt,a4paper]{ctexart}
\usepackage{amsmath,amssymb,geometry,enumitem,hyperref}
\geometry{margin=2.1cm}\setlist{itemsep=2pt,topsep=3pt}
\hypersetup{colorlinks=true,linkcolor=blue}
\newcommand{\EN}[1]{\par\noindent \textbf{English.} #1\par}
\newcommand{\READ}[1]{\par\noindent \textbf{公式英文读法 / How to read it.} #1\par}
\newcommand{\EX}[2]{\par\noindent \textbf{例子 / Example.} #1\par\EN{#2}}
\newcommand{\ASK}[4]{\par\noindent \textbf{可能问题 / Possible question.} #1\par\EN{#2}\noindent \textbf{回答要点 / Key answer.} #3\par\EN{#4}\medskip}
\title{代数拓扑 APS 面试准备\\[4pt]\large Croom《Basic Concepts of Algebraic Topology》第 1--5 章}
\author{Algebraic Topology: Bilingual Review}
\date{}
\begin{document}\maketitle
\section*{课程概述 | Course overview}
\noindent \textbf{中文口述稿。} 在代数拓扑课上，我学习用代数方法区分形状。比如一个三角形边框由三个点和三条边组成，沿边走一圈会回到起点；如果没有填入三角面，这个环不能被已有的面填满。它的一维同调群是 $H_1\cong\mathbb Z$，整数可以理解为绕环的次数。填入三角面后，这个环成为面的边界，一维同调变为零。这个例子说明同调怎样把几何中的“洞”变成可以计算的代数信息；它也联系了群论和几何。
\EN{In algebraic topology, I learned to distinguish shapes using algebra. A triangular frame has three vertices and three edges. Following the edges once gives a closed loop, and without a triangular face that loop cannot be filled by the available simplices. Its first homology group is $H_1\cong\mathbb Z$, where integers correspond to multiples of the loop. After filling the triangle, the loop becomes the boundary of the face and the first homology is zero. This shows how homology turns a geometric hole into an algebraic object that can be computed, connecting group theory and geometry.}

\noindent \textbf{具体应用联系。} 拓扑数据分析把传感器点云、三维扫描或分子构形近似成 simplicial complex，再计算连通分支和环洞。若数据来自一个圆环形状，理想情况下应有一个连通分支和一个一维洞；测量噪声会制造很小、很短暂的假洞，因此实践中常比较不同尺度下拓扑特征是否持续存在。同调在同伦等价下不变；任意连续映射却可能压缩掉一个环，所以不能说任何连续变形都保持洞的数量。
\EN{Topological data analysis approximates sensor point clouds, 3D scans, or molecular configurations by a simplicial complex and computes components and loops. An ideal circle has one connected component and one one-dimensional hole. Noise can create small, short-lived false holes, so practical methods compare features across scales. Homology is invariant under homotopy equivalence, but an arbitrary continuous map can collapse a loop and need not preserve holes.}

\section{几何复形与方向 | Geometric complexes and orientation}
\noindent $k$-单形是 $k+1$ 个仿射独立点的凸包：$0$-单形是顶点，$1$-单形是边，$2$-单形是三角形，$3$-单形是四面体。单纯复形由一组单形组成，并满足任意面的面仍在其中、两个单形的交是共同的面或空集。几何复形的并称多面体。给边或三角形定向，是为了在边界公式中保留符号。
\EN{A $k$-simplex is the convex hull of $k+1$ affinely independent points: vertices, edges, triangles, and tetrahedra are the cases $k=0,1,2,3$. A simplicial complex contains the faces of each simplex and has intersections that are common faces or empty. Orientations assign signs so that boundary terms can cancel correctly.}

\section{链、圈、边界与同调 | Chains, cycles, boundaries, and homology}
\noindent 在系数域或整数环 $R$ 上，$C_k(K)$ 是由有向 $k$-单形生成的自由模，元素称 $k$-链。边界算子 $\partial_k:C_k\to C_{k-1}$ 对有向单形定义为
\[
\partial_k[v_0,\ldots,v_k]=\sum_{j=0}^k(-1)^j[v_0,\ldots,\widehat v_j,\ldots,v_k].
\]
帽号表示删去该顶点。每个面以交替符号出现，故 $\partial_{k-1}\partial_k=0$：边界的边界为零。圈是满足 $\partial_k c=0$ 的链；边界是形如 $\partial_{k+1}b$ 的链。因为边界都是圈，商群
\[
H_k(K)=\ker\partial_k/\operatorname{im}\partial_{k+1}
\]
有定义。它把“没有边界的闭合对象”按“本身是某个更高维对象的边界”视为平凡。
\EN{With integer or field coefficients, $C_k(K)$ consists of finite linear combinations of oriented $k$-simplices. The boundary map is the alternating sum shown above; the hat means that a vertex is omitted. Applying the boundary twice gives zero. A cycle lies in the kernel, while a boundary lies in the image of the next boundary map. Every boundary is a cycle, so the homology group is the quotient of cycles by boundaries.}
\READ{“H k is the kernel of the boundary map in dimension k, modulo the image of the boundary map in dimension k plus one.”}

\EX{把三角形边界定向为 $[v_0,v_1]+[v_1,v_2]+[v_2,v_0]$。三个顶点的边界项成对抵消，所以该链是 $1$-圈；若填入三角形面，它正是该 $2$-单形的边界，因而在同调群中代表零。圆周只由边组成而没有填充面，其绕一圈的 $1$-圈不是边界，因此保留一个洞。}{Orient the boundary of a triangle as $[v_0,v_1]+[v_1,v_2]+[v_2,v_0]$. Its vertex terms cancel, so it is a 1-cycle. If the triangular face is included, that cycle is a boundary and represents zero in homology. A circle has no filling face, so its loop is not a boundary and detects a hole.}

\noindent\textbf{完整计算：三角形边界。} 以整数为系数，令 $e_{01}=[v_0,v_1]$、$e_{12}=[v_1,v_2]$、$e_{20}=[v_2,v_0]$，则 $\partial e_{01}=v_1-v_0$，其余两条边同理。对任意 $1$-链 $c=ae_{01}+be_{12}+ce_{20}$，令 $\partial c=0$ 等价于 $a=b=c$，故 $\ker\partial_1$ 由 $e_{01}+e_{12}+e_{20}$ 生成。由于没有 $2$-单形，$\operatorname{im}\partial_2=0$，因此 $H_1\cong\mathbb Z$。三个顶点连在一起，所以 $H_0\cong\mathbb Z$。填入三角形后这个圈成为面的边界，$H_1$ 变为零。
\EN{For the boundary of a triangle, orient the edges cyclically. An arbitrary 1-chain has zero boundary exactly when all three edge coefficients are equal, so the cycle group is generated by the loop around the triangle. There is no 2-simplex, so the boundary group in dimension one is zero and $H_1\cong\mathbb Z$. All vertices are connected, so $H_0\cong\mathbb Z$. Filling the triangle makes the loop a boundary and changes $H_1$ to zero.}

\noindent\textbf{Betti 数。} 当使用域系数时，$\beta_k=\dim H_k$ 记录第 $k$ 维独立同调类个数：$\beta_0$ 是连通分支数，$\beta_1$ 常解释为独立环洞数，$\beta_2$ 在常见三维模型中可检测被包围的空腔。解释依赖空间及系数；整数同调还可能有有限阶挠子群，单靠 Betti 数看不到。
\EN{With field coefficients, the Betti number $\beta_k=\dim H_k$ counts independent homology classes in dimension $k$. The zeroth counts connected components; the first often counts independent loops; the second can detect enclosed voids in familiar three-dimensional spaces. The interpretation depends on the space and coefficients. Integer homology can also contain torsion that Betti numbers do not detect.}

\noindent 有限单纯复形的 Euler 特征为 $\chi(K)=\sum_k(-1)^k f_k$，其中 $f_k$ 是 $k$-单形个数；Euler--Poincaré 公式给出 $\chi(K)=\sum_k(-1)^k\operatorname{rank}H_k(K)$（使用合适的系数并计入 Betti 数）。它让容易数出的顶点、边、面与拓扑洞数产生联系，但 Euler 特征相同不保证空间同胚。
\EN{For a finite simplicial complex, the Euler characteristic is the alternating sum of the numbers $f_k$ of simplices. The Euler--Poincare formula expresses it as the alternating sum of the ranks of the homology groups, or Betti numbers. It connects a combinatorial count of vertices, edges, and faces to topological holes, although equal Euler characteristics do not imply that spaces are homeomorphic.}
\EX{四面体表面有 $f_0=4,f_1=6,f_2=4$，所以 $\chi=4-6+4=2$，与球面相同。环面有 $\chi=0$，并有两个独立的一维基本环方向。}{The boundary of a tetrahedron has four vertices, six edges, and four faces, so $\chi=4-6+4=2$, as for a sphere. A torus has Euler characteristic zero and two independent one-dimensional loop directions.}

\section{单纯映射与基本群 | Simplicial maps and the fundamental group}
\noindent 单纯映射把顶点送到顶点，并把每个单形送到某个单形（允许维数下降），因此诱导链映射并保持同调。连续映射常可在足够细分后由单纯映射近似，这就是单纯逼近思想：先把连续几何问题转成有限组合数据，再证明代数不变量不依赖所选细分。
\EN{A simplicial map sends vertices to vertices and each simplex to a simplex, possibly of lower dimension. It induces maps on chains and homology. After sufficient subdivision, continuous maps can often be approximated by simplicial maps, converting geometric questions into finite combinatorial ones.}

\noindent 以点 $x_0$ 为基点，基于 $x_0$ 的闭路按保持基点的同伦分类；闭路同伦类在路径拼接下构成基本群 $\pi_1(X,x_0)$。它记录空间中不能连续缩成常值路的环。圆周的基本群同构于 $\mathbb Z$，整数表示绕圈次数；球面 $S^2$ 的基本群平凡。若空间路径连通，改变基点会得到同构群，但同构通常依赖选取的路径。
\READ{“The fundamental group consists of based loops modulo based homotopy, with path concatenation as its operation.”}

\section{覆盖空间 | Covering spaces}
\noindent 映射 $p:\widetilde X\to X$ 是覆盖映射，若每个 $x\in X$ 都有邻域 $U$，使 $p^{-1}(U)$ 是若干互不相交开集的并，且每一片在 $p$ 下同胚到 $U$。局部看，覆盖空间与底空间相同；全局上，一个底空间点可以有多个原像。路径提升与同伦提升把底空间中的闭路信息转成覆盖空间中的端点变化，因此覆盖空间与基本群的子群结构相连。
\EN{A covering map is locally a disjoint union of copies: each point has a neighborhood whose inverse image consists of disjoint open sets mapped homeomorphically onto it. Locally the covering space looks like the base, while globally a point can have several preimages. Path and homotopy lifting connect loops in the base to endpoint changes upstairs and relate coverings to subgroups of the fundamental group.}

\EX{指数映射 $p:\mathbb R\to S^1$，$p(t)=e^{2\pi it}$，是实直线对圆周的覆盖。圆周上绕一圈的闭路从 $t=0$ 提升到实线后终点为 $1$，没有回到起点；绕 $n$ 圈则终点为 $n$。因此提升记录了绕数，并显示 $\pi_1(S^1)\cong\mathbb Z$。}{The map $p:\mathbb R\to S^1$ given by $p(t)=e^{2\pi it}$ covers the circle. A loop going once around the circle lifts from $0$ to $1$ in the real line, so its lift does not return to its starting point. A loop going around $n$ times ends at $n$. This records the winding number and explains why $\pi_1(S^1)\cong\mathbb Z$.}

\noindent\textbf{路径提升如何用。} 固定提升起点 $\widetilde x_0$ 后，底空间的每条路径都有唯一提升。若底空间的两个基点闭路保持端点同伦，则它们从同一提升起点出发的提升终点相同。因此圆周中绕一圈和常值路不可能基点同伦：前者提升终点是 $1$，后者是 $0$。
\EN{After fixing a starting point in the covering space, a path in the base has a unique lift. Based-homotopic loops have lifts with the same endpoint when started at the same point. On the circle, a once-around loop and the constant loop therefore cannot be based-homotopic: their lifts from zero end at one and zero, respectively.}

\ASK{为什么要把圈除以边界来定义同调群？}{Why is homology defined as cycles modulo boundaries?}{圈满足边界为零，表示闭合结构；若它是某个高一维链的边界，就视为可填充。取 $\ker\partial_k/\operatorname{im}\partial_{k+1}$，并把相差边界的圈视为同一个同调类。}{Cycles are closed chains. A cycle that bounds a chain one dimension higher is fillable and represents zero. The quotient identifies cycles whose difference is a boundary.}
\ASK{圆周的基本群为什么是整数群？}{Why is the fundamental group of the circle the integers?}{闭路可按绕圆周的有向次数分类，绕数是整数；路径拼接对应绕数相加，零次绕数的闭路可缩为一点。因此群运算对应整数加法。}{Loops are classified by their signed winding number, an integer. Concatenation adds winding numbers, and a loop with winding number zero is contractible, so the group operation is integer addition.}
\ASK{同调群与基本群有什么区别？}{How do homology and the fundamental group differ?}{基本群记录基点闭路及其可能非交换的拼接结构；同调群把链线性化并取商，通常更容易计算，但会丢失部分非交换信息。两者在同伦等价下不变；基本群还要适当处理基点。}{The fundamental group records based loops with a possibly noncommutative concatenation law. Homology linearizes chains and takes a quotient, losing some noncommutative information. Both are invariant under homotopy equivalence, with appropriate basepoints for the fundamental group.}
\section{从零开始：同胚、同伦与拓扑不变量 | Homeomorphism, homotopy, and invariants}
\noindent \textbf{同胚。} 若 $f:X\to Y$ 是连续双射且逆映射也连续，则 $X,Y$ 同胚。它允许拉伸和弯曲，但不能撕开或粘接；例如圆周与正方形边界同胚。仅有连续双射一般不够，还需要逆连续。
\EN{A homeomorphism is a continuous bijection with a continuous inverse. It permits bending and stretching but preserves the topological structure. For example, a circle is homeomorphic to the boundary of a square. A continuous bijection alone need not be a homeomorphism.}
\noindent \textbf{同伦。} 两个连续映射 $f,g:X\to Y$ 若存在连续 $H:X\times[0,1]\to Y$ 满足 $H(x,0)=f(x)$、$H(x,1)=g(x)$，则它们同伦。$t$ 是连续变形参数，不必解释为物理时间。若存在 $f:X\to Y$、$g:Y\to X$，使 $gf\simeq\mathrm{id}_X$、$fg\simeq\mathrm{id}_Y$，则 $X,Y$ 同伦等价。同胚蕴含同伦等价，反之不一定。
\EN{Maps $f,g:X\to Y$ are homotopic if a continuous $H:X\times[0,1]\to Y$ deforms one into the other. Spaces are homotopy equivalent if maps in both directions compose to maps homotopic to the respective identities. Homeomorphism implies homotopy equivalence, but not conversely.}
\EX{实直线 $\mathbb R$ 与一点不是同胚：它们的点数不同；但 $H(x,t)=(1-t)x$ 把恒等映射连续收缩到常值 $0$，所以 $\mathbb R$ 与一点同伦等价。圆周不能这样收缩，因为绕圆一圈的闭路有非零绕数。}{The real line is not homeomorphic to a point, but $H(x,t)=(1-t)x$ contracts it to zero, making it homotopy equivalent to a point. A circle cannot be contracted in the same way because a loop around it has nonzero winding number.}
\ASK{为什么比较基本群能证明两个空间不等价？}{How can fundamental groups distinguish spaces?}{同伦等价的空间有同构的基本群（基点按路径相应处理）。若基本群不同构，则它们不同伦等价，更不可能同胚。例如一点的基本群平凡，而圆周的基本群是 $\mathbb Z$。}{Homotopy-equivalent spaces have isomorphic fundamental groups, with compatible basepoints. If the groups are not isomorphic, the spaces are not homotopy equivalent and cannot be homeomorphic. A point has trivial fundamental group, while a circle has fundamental group $\mathbb Z$.}

\section{单形、链与边界：逐项计算 | Simplices, chains, and boundaries}
\noindent 有向边 $[v_i,v_j]$ 的边界是终点减起点：$\partial_1[v_i,v_j]=[v_j]-[v_i]$。有向三角形的边界为
\[\partial_2[v_0,v_1,v_2]=[v_1,v_2]-[v_0,v_2]+[v_0,v_1].\]
改变单形的方向会改变符号。对上式再取边界，每个顶点出现两次且符号相反，因此 $\partial_1\partial_2=0$。推广到任意维数，每个低两维的面因删去两个顶点的先后次序不同出现两次，符号相反。
\EN{The boundary of an oriented edge is endpoint minus starting point. The boundary of an oriented triangle is the alternating sum shown above. Applying the boundary operator again makes every vertex cancel. In all dimensions, every face of codimension two occurs twice with opposite signs, giving $\partial^2=0$.}
\READ{“The boundary of the oriented triangle v zero, v one, v two is the edge v one to v two minus the edge v zero to v two plus the edge v zero to v one.”}
\EX{只有一条边 $[v_0,v_1]$ 的复形中，$C_1\cong\mathbb Z$、$C_0\cong\mathbb Z^2$。$\partial_1(a[v_0,v_1])=a(v_1-v_0)$，所以 $\ker\partial_1=0$，$H_1=0$。在 $H_0$ 中商去 $v_1-v_0$ 后两个顶点代表同一个连通分支，故 $H_0\cong\mathbb Z$。}{For a complex with one edge and its two endpoints, $C_1\cong\mathbb Z$ and $C_0\cong\mathbb Z^2$. The boundary of $a[v_0,v_1]$ is $a(v_1-v_0)$, so $\ker\partial_1=0$ and $H_1=0$. Quotienting $C_0$ by the image identifies the endpoints, leaving $H_0\cong\mathbb Z$ for one connected component.}

\section{算同调时的五个步骤 | Five steps for computing homology}
\noindent 实际计算可按以下顺序：列出各维有向单形及链群 $C_k$；写出每个边界算子 $\partial_k$；解线性方程求圈群 $Z_k=\ker\partial_k$；由高一维单形求边界群 $B_k=\operatorname{im}\partial_{k+1}$；最后计算商 $H_k=Z_k/B_k$。用域系数时，可对边界矩阵做行变换；用整数系数时，还要考虑挠，不能仅靠普通实数矩阵的秩。
\EN{To compute homology, list the oriented simplices and chain groups, write the boundary maps, find cycles as kernels, find boundaries as images from the next dimension, and form the quotient. Over a field, row reduction gives dimensions. With integer coefficients, torsion must also be considered, so real-matrix ranks alone are insufficient.}
\EX{系数取 $\mathbb R$，三角形边界有三个顶点和三条边、没有面。$\dim C_1=3$，边界矩阵的秩为 $2$，所以 $\dim Z_1=3-2=1$；由于 $C_2=0$，有 $B_1=0$，得 $\beta_1=1$。同时 $\beta_0=3-2=1$。若填入一个三角面，则 $\partial_2$ 的像正好张成原有那个圈，$\beta_1$ 变为 $0$。}{Over $\mathbb R$, the triangular boundary has three vertices and three edges and no faces. The edge-boundary matrix has rank two, so its kernel has dimension $3-2=1$. Since $C_2=0$, $\beta_1=1$, and $\beta_0=3-2=1$. Adding the triangular face makes the cycle a boundary, so $\beta_1$ becomes zero.}

\section{基本群：闭路的组合与绕数 | Fundamental group and winding number}
\noindent 以 $x_0$ 为基点的闭路是连续映射 $\gamma:[0,1]\to X$，且 $\gamma(0)=\gamma(1)=x_0$。两条闭路可先走第一条再走第二条；取保持端点的同伦等价类后，拼接在同伦类层面满足结合律。常值路是单位元，反向走的路径给逆元。群可能非交换，取决于空间结构。
\EN{A based loop is a continuous map $\gamma:[0,1]\to X$ with both endpoints at $x_0$. Concatenating loops gives the group operation on based-homotopy classes. The constant loop is the identity and the reversed loop gives an inverse. The group can be noncommutative.}
\EX{在 $S^1$ 上取 $\gamma_n(t)=e^{2\pi int}$。$n=2$ 表示逆时针绕两圈，$n=-1$ 表示顺时针绕一圈。先走 $\gamma_2$ 再走 $\gamma_{-1}$，总绕数为 $1$，所以其同伦类对应整数 $1$。通过 $\mathbb R\to S^1$ 的路径提升可证明绕数在基点同伦下不变。}{On $S^1$, $\gamma_n(t)=e^{2\pi int}$ winds $n$ times. The loop with $n=2$ followed by the loop with $n=-1$ has total winding one. Lifting paths to $\mathbb R$ proves that winding number is unchanged under based homotopy.}

\section{覆盖空间：从局部复制到全局信息 | Coverings and lifting}
\noindent 覆盖映射的每个底空间点附近都有\emph{均匀覆盖邻域} $U$：$p^{-1}(U)$ 分成互不相交的开集，每一片都通过 $p$ 同胚到 $U$。取圆周上不绕满一整圈的小弧，它在实线覆盖 $t\mapsto e^{2\pi it}$ 下的原像是相隔整数距离的一列小区间。一个点上方的原像集合称纤维。
\EN{An evenly covered neighborhood $U$ has inverse image equal to a disjoint union of open sets, each mapped homeomorphically onto $U$. For the covering $t\mapsto e^{2\pi it}$ of the circle, a small arc has inverse-image intervals separated by integers. The inverse image of a single point is called its fiber.}
\EX{映射 $p:S^1\to S^1$、$p(z)=z^2$ 是双层覆盖：对任意点 $w\in S^1$，其原像恰有两个平方根。底空间沿圆周绕一圈的闭路提升到上层圆周后，会从一个平方根走到另一个；绕两圈才回到原来的提升起点。}{The map $p:S^1\to S^1$, $p(z)=z^2$, is a two-sheeted covering: each $w\in S^1$ has two square roots on the circle. A loop going once around the base lifts to a path from one root to the other. Going around twice returns to the starting lift.}
\ASK{为什么覆盖空间能帮助计算基本群？}{Why do coverings help compute fundamental groups?}{固定提升起点后，闭路的提升终点记录绕行信息；基点同伦保持该终点。覆盖空间可能更简单，例如实线可收缩，而它覆盖的圆周有非平凡闭路。}{With a fixed lift starting point, the endpoint of a lifted loop records winding information and is unchanged by based homotopy. The covering space can be simpler: the real line is contractible even though its base circle has nontrivial loops.}
\section{零基础预备：集合、映射、群与商 | Prerequisites}
\noindent \textbf{集合与映射。} 集合是对象的整体，$x\in X$ 表示 $x$ 属于 $X$。映射 $f:X\to Y$ 给每个 $x\in X$ 指定唯一一个 $f(x)\in Y$；$f$ 单射表示不同输入不会有同一输出，满射表示 $Y$ 中每个点都是某个输入的输出。双射同时满足两者。
\EN{A set is a collection of objects, and $x\in X$ means that $x$ belongs to $X$. A map $f:X\to Y$ assigns exactly one output in $Y$ to every input in $X$. It is injective if distinct inputs have distinct outputs, surjective if every element of $Y$ is reached, and bijective if both hold.}
\EX{令 $X=\{1,2,3\}$、$Y=\{a,b\}$，取 $f(1)=a$、$f(2)=a$、$f(3)=b$。该映射满射，因为 $a,b$ 都被取到；它不是单射，因为 $1,2$ 的输出相同。}{Let $X=\{1,2,3\}$ and $Y=\{a,b\}$, with $f(1)=a$, $f(2)=a$, and $f(3)=b$. The map is surjective because both elements of $Y$ are reached, but it is not injective because one and two have the same output.}

\noindent \textbf{群。} 群 $G$ 是带有运算的集合，该运算满足结合律、存在单位元、每个元素有逆元。若还满足交换律，称阿贝尔群。群同态 $h:G\to H$ 保持运算：$h(ab)=h(a)h(b)$；核是映到单位元的元素，像是实际取到的元素。同构是双射同态，表示两个群的运算结构相同。
\EN{A group is a set with an associative operation, an identity element, and inverses. If its operation is commutative, it is abelian. A group homomorphism preserves the operation. Its kernel consists of elements mapped to the identity, and its image consists of all outputs. A bijective homomorphism is an isomorphism.}
\EX{整数集合 $\mathbb Z$ 在加法下是阿贝尔群：单位元是 $0$，$n$ 的逆元是 $-n$。映射 $h:\mathbb Z\to\mathbb Z_2$ 定义为 $h(n)=n\bmod2$，它把加法映成模 $2$ 加法，核为偶整数 $2\mathbb Z$，像为 $\mathbb Z_2$。}{The integers under addition form an abelian group: the identity is zero and the inverse of $n$ is $-n$. The map $h:\mathbb Z\to\mathbb Z_2$ sending an integer to its parity is a homomorphism. Its kernel is the even integers $2\mathbb Z$, and its image is all of $\mathbb Z_2$.}

\noindent \textbf{等价关系与商。} 等价关系是满足自反、对称和传递性的关系。$x$ 的等价类 $[x]$ 是所有与 $x$ 等价的对象；商集由这些等价类组成。在阿贝尔群 $G$ 中，子群 $H$ 给出关系 $a\sim b\Leftrightarrow a-b\in H$，商群 $G/H$ 的元素是陪集 $a+H$。
\EN{An equivalence relation is reflexive, symmetric, and transitive. The equivalence class $[x]$ contains all elements equivalent to $x$, and the quotient set consists of these classes. For an abelian group $G$ and subgroup $H$, define $a\sim b$ when $a-b\in H$; the quotient group consists of cosets $a+H$.}
\EX{在 $\mathbb Z$ 中按相差偶数分类，得到两类：$0+2\mathbb Z$ 是所有偶数，$1+2\mathbb Z$ 是所有奇数。因此 $\mathbb Z/2\mathbb Z$ 有两个元素，奇数类加奇数类得到偶数类。同调群中的“圈模去边界”使用相同的商思想。}{Classify integers by whether their difference is even. The two classes are $0+2\mathbb Z$, the even integers, and $1+2\mathbb Z$, the odd integers. Their quotient has two elements, and odd plus odd is even. Homology uses the same quotient idea with cycles modulo boundaries.}

\section{几何单形与复形：逐个画出来 | Building a simplicial complex}
\noindent \textbf{单形。} 从几何上看，零维单形是一个点 $[v_0]$；一维单形是两点间整条线段 $[v_0,v_1]$，包含端点；二维单形是实心三角形 $[v_0,v_1,v_2]$，包含三条边及其端点；三维单形是实心四面体。其维数比顶点数少一。面由选取部分顶点得到。
\EN{A zero-simplex is a vertex, a one-simplex is a filled line segment including its endpoints, a two-simplex is a filled triangle including its edges, and a three-simplex is a filled tetrahedron. Its dimension is one less than its number of vertices. Faces are obtained by selecting subsets of vertices.}
\EX{设 $v_0=(0,0)$、$v_1=(1,0)$、$v_2=(0,1)$。二维单形是 $\{(x,y):x\ge0,y\ge0,x+y\le1\}$；它的三个一维面分别是 $x=0$、$y=0$、$x+y=1$ 上的相应线段。只取三条边得到三角形边界，加入内部才得到实心三角形。}{For vertices $(0,0)$, $(1,0)$, and $(0,1)$, the filled triangle is $\{(x,y):x\ge0,y\ge0,x+y\le1\}$. Its three one-dimensional faces are its edges. Keeping only the edges gives the boundary; including the interior gives the filled triangle.}

\noindent \textbf{单纯复形。} 一组单形构成单纯复形，要求每个单形的每个面也在复形中，且两个单形的交只能是共同面或空集。记它的几何实现为 $|K|$，即把这些单形视为几何图形后的并。该条件防止两个三角形以不沿公共面的方式交叉。
\EN{A simplicial complex is a collection of simplices containing every face of each simplex, where any two simplices intersect only in a common face or not at all. Its geometric realization $|K|$ is the union of the geometric simplices.}
\EX{复形 $K=\{[v_0],[v_1],[v_2],[v_0,v_1],[v_1,v_2],[v_2,v_0]\}$ 包含三个顶点和三条边，不包含三角面，其几何实现是一个环。若把 $[v_0,v_1,v_2]$ 加入，原有各面都已经在 $K$ 中，得到实心三角形。}{The complex consisting of three vertices and their three edges, without a triangular face, has a loop as its realization. Adding the two-simplex gives a filled triangle; all its faces are already present.}

\section{链群到同调群：每一步代表什么 | From chain groups to homology}
\noindent \textbf{链群。} $C_k(K;\mathbb Z)$ 是所有有向 $k$-单形的有限整数线性组合。系数 $2$ 表示该单形取两次，系数 $-1$ 表示反向取一次，系数 $0$ 表示不使用；相加时相同单形的系数相加。它是一个阿贝尔群。
\EN{The chain group $C_k(K;\mathbb Z)$ consists of finite integer linear combinations of oriented $k$-simplices. A coefficient two means two copies, minus one reverses orientation, and zero omits the simplex. Chains add by adding coefficients and form an abelian group.}
\EX{如果 $K$ 的边只有 $e_1,e_2$，则 $C_1(K;\mathbb Z)=\{ae_1+be_2:a,b\in\mathbb Z\}\cong\mathbb Z^2$。例如 $2e_1-e_2$ 是一条 $1$-链；它是否是圈还要计算边界。}{If the only edges are $e_1,e_2$, then $C_1(K;\mathbb Z)=\{ae_1+be_2:a,b\in\mathbb Z\}\cong\mathbb Z^2$. The chain $2e_1-e_2$ is a valid one-chain; whether it is a cycle requires computing its boundary.}

\noindent \textbf{边界算子。} $\partial_k:C_k\to C_{k-1}$ 先对每个有向单形取带符号的面，再按线性性扩展。$\partial^2=0$ 使每个边界自动是圈，即 $B_k=\operatorname{im}\partial_{k+1}\subseteq Z_k=\ker\partial_k$。
\EN{The boundary map $\partial_k:C_k\to C_{k-1}$ takes the signed faces of each simplex and extends linearly. Because $\partial^2=0$, every boundary is automatically a cycle: $B_k=\operatorname{im}\partial_{k+1}$ is contained in $Z_k=\ker\partial_k$.}
\EX{三条定向边 $e_{01}=[v_0,v_1]$、$e_{12}=[v_1,v_2]$、$e_{20}=[v_2,v_0]$ 的和 $z=e_{01}+e_{12}+e_{20}$ 满足 $\partial z=(v_1-v_0)+(v_2-v_1)+(v_0-v_2)=0$，所以 $z$ 是圈。如果复形包含三角面，则 $z=\partial[v_0,v_1,v_2]$，它同时是边界。}{For cyclic edges $e_{01},e_{12},e_{20}$, the chain $z=e_{01}+e_{12}+e_{20}$ has boundary $(v_1-v_0)+(v_2-v_1)+(v_0-v_2)=0$, so it is a cycle. If the triangular face is present, it is also the boundary of that face.}

\noindent \textbf{同调类。} 两个 $k$-圈 $z,z'$ 若相差一个边界，即 $z-z'=\partial_{k+1}c$，就视为同一个同调类；$H_k=Z_k/B_k$。零类包括所有边界。直观地，同调忽略可以由更高维单形填掉的差别，保留不能这样填掉的闭合结构。
\EN{Two $k$-cycles represent the same homology class if their difference is a boundary. The homology group is $H_k=Z_k/B_k$, and its zero class consists of the boundaries. Intuitively, homology ignores changes that can be filled by higher-dimensional simplices.}
\EX{只含三角形边界时，绕环一次的圈 $z$ 不是边界，故 $[z]\ne0$；$2z$ 绕两次，代表 $2[z]$。若加入三角面，$z$ 成为边界，所以 $[z]=0$，并且 $[2z]=0$。两个空间仅差一个面，却有不同的一维同调。}{For the unfilled triangular boundary, the once-around cycle $z$ is not a boundary, so $[z]\ne0$, and twice around gives $2[z]$. After adding the face, $z$ is a boundary, so both $[z]$ and $[2z]$ are zero. Adding one face changes the first homology.}

\section{典型空间同调：怎样回答一个洞 | Homology of familiar spaces}
\noindent \textbf{连通分支与 $H_0$。} 对非空有限单纯复形，整数系数 $H_0\cong\mathbb Z^c$，其中 $c$ 是连通分支数。每个分支贡献一个独立生成元：边的边界把同一分支内的顶点识别起来，不会把两个互不相连的分支识别起来。
\EN{For a nonempty finite simplicial complex, $H_0\cong\mathbb Z^c$, where $c$ is the number of connected components. Edge boundaries identify vertices in the same component, while vertices in separate components remain independent.}
\EX{两个互不相连的点没有边，故 $C_0\cong\mathbb Z^2$、$\operatorname{im}\partial_1=0$，因此 $H_0\cong\mathbb Z^2$。如果加一条连接它们的边，则边界是两点之差，商掉这条关系后 $H_0\cong\mathbb Z$。}{Two isolated points have $C_0\cong\mathbb Z^2$ and no edge boundaries, so $H_0\cong\mathbb Z^2$. Adding an edge makes the difference of the two vertices a boundary, leaving $H_0\cong\mathbb Z$.}

\noindent \textbf{圆周、圆盘和球面。} 对整数系数，圆周有 $H_0(S^1)\cong\mathbb Z$、$H_1(S^1)\cong\mathbb Z$，更高维同调为零；实心圆盘可收缩，有 $H_0\cong\mathbb Z$、$H_k=0$（$k\ge1$）；二维球面有 $H_0(S^2)\cong\mathbb Z$、$H_2(S^2)\cong\mathbb Z$、$H_1(S^2)=0$。球面的一维圈可以在球面上缩到一点，但球面整体给出一个非平凡二维类。
\EN{With integer coefficients, the circle has $H_0\cong\mathbb Z$ and $H_1\cong\mathbb Z$. A filled disk is contractible, so its positive-dimensional homology is zero. The sphere $S^2$ has $H_0\cong\mathbb Z$, $H_2\cong\mathbb Z$, and $H_1=0$. Loops on a sphere contract, but the whole oriented sphere represents a two-dimensional class.}
\EX{把四面体表面看作球面的三角剖分：顶点数 $4$、边数 $6$、面数 $4$，Euler 特征 $4-6+4=2$。它的一维洞数为 $0$，二维基本类为 $1$，因此 Betti 数满足 $\beta_0-\beta_1+\beta_2=1-0+1=2$，与单形计数吻合。}{The tetrahedron boundary triangulates a sphere, with four vertices, six edges, and four faces. Its Euler characteristic is $4-6+4=2$. Its Betti numbers give $\beta_0-\beta_1+\beta_2=1-0+1=2$, agreeing with the simplex count.}

\section{单纯映射与连续变形：一个具体例子 | Simplicial maps}
\noindent \textbf{单纯映射定义。} 给定复形 $K,L$，顶点映射 $f$ 若把 $K$ 中每个单形的顶点集合映到 $L$ 中某个单形的顶点集合，就称单纯映射；不同顶点允许映到同一顶点，因此维数可以下降。它在线性延拓后给出连续映射，并诱导同调群之间的同态。
\EN{A simplicial map sends vertices of $K$ to vertices of $L$ so that the vertices of every simplex map into a simplex of $L$. Vertices may be identified, reducing dimension. Linear extension gives a continuous map and induces homomorphisms on homology.}
\EX{把三角形边界的三个顶点都映到单点复形中的唯一顶点，三条边都被压成这个点。原先绕一圈的非零 $H_1$ 类被送到单点的 $H_1=0$。这具体展示了连续映射如何诱导同调映射，也展示映射不一定保留所有洞。}{Map every vertex of the triangular boundary to the unique vertex of a one-point complex. All edges collapse to the point, and the nonzero loop class maps to zero in the point's $H_1$. An induced map on homology need not preserve every hole.}

\section{面试追问：先说定义，再举例 | Oral questions}
\ASK{圈一定是边界吗？}{Is every cycle a boundary?}{不一定。边界一定是圈，因为 $\partial^2=0$；但只含边的三角形边界中，绕一圈是圈，由于没有三角面，它不是边界。填入三角面后它才成为边界。}{No. Every boundary is a cycle because $\partial^2=0$, but the loop around an unfilled triangle is a cycle and is not a boundary. Adding the triangular face makes it a boundary.}
\ASK{圆周和圆盘都是连通的，为什么不同胚？}{Both a circle and a disk are connected. Why are they not homeomorphic?}{它们的 $H_0$ 都是 $\mathbb Z$，但圆周的 $H_1\cong\mathbb Z$，圆盘的 $H_1=0$。同胚会诱导同构的同调群，所以它们不同胚。}{Their $H_0$ groups are both $\mathbb Z$, but the circle has $H_1\cong\mathbb Z$ and the disk has $H_1=0$. Homeomorphic spaces have isomorphic homology groups, so these spaces are not homeomorphic.}
\ASK{如何向非数学背景的人解释同调？}{How would you explain homology to someone without a mathematics background?}{先把形状拆成点、边和面，再寻找闭合的边或面。若一个闭合边界可以用已有的面填满，就不把它算成洞；若不能填满，它会留下一个可计算的同调类。三角形边框与填满的三角形是最简单的对照。}{First break a shape into vertices, edges, and faces, then look for closed loops or surfaces. A loop that is filled by available faces does not count as a hole; an unfilled loop can represent a nonzero homology class. An empty triangular frame and a filled triangle are simple contrasting examples.}
\end{document}
